% --------------------------------------------------------------
% Small Integrated Network Project Report
% --------------------------------------------------------------

\documentclass[11pt,a4paper]{article}
\usepackage[margin=1in]{geometry}
\usepackage[hidelinks]{hyperref}
\usepackage{graphicx}
\usepackage{float}
\usepackage{diagbox}
\usepackage{booktabs}
\usepackage{tabularx}
\usepackage{array}
\usepackage{amsmath,amssymb}
\usepackage{enumitem}
\usepackage{pdflscape}

\title{Small Integrated Network Project Report}
\author{}
\date{\today}

\begin{document}

\maketitle

\section{Problem Statement}
% TODO: Define the core problem or objective of your small integrated network project.
% TODO: Specify the domain chosen (e.g., Hospital, University, Smart Home, or Accounts).

% Example placeholder:
% This project addresses the need for secure, automated network segmentation in a \textbf{[Hospital/University/Smart Home/Accounts]} environment. The core objective is to design and implement a network infrastructure that enables reliable communication between designated zones while enforcing security policies and preventing unauthorized access. Specific challenges include [briefly describe the problem, e.g., "managing traffic between patient wards and administrative offices" / "connecting dormitory networks to campus resources" / "integrating IoT devices in a smart home" / "securing financial data in an accounts system"].

\section{Requirements}
% TODO: Hardware, software, and technology requirements needed for the network setup.
% TODO: List the maximum technologies covered in class that are integrated into this project.

% Hardware Requirements:
\begin{itemize}
  \item [Computers/Devices:] \textit{[e.g., 2 routers, 1 switch, 10 end-devices]}
  \item [Network Interfaces:] \textit{[e.g., Gigabit Ethernet ports, Wi-Fi adapters]}
  \item [Other Hardware:] \textit{[e.g., UPS, access points, IoT sensors]}
\end{itemize}

% Software Requirements:
\begin{itemize}
  \item [Operating Systems:] \textit{[e.g., Ubuntu 22.04, Windows 10, Raspberry Pi OS]}
  \item [Network Software:] \textit{[e.g., Cisco Packet Tracer, GNS3, Mininet, Wireshark]}
  \item [Programming Languages/Tools:] \textit{[e.g., Python 3, Socket programming, simulation frameworks]}
\end{itemize}

% Technologies Covered (maximum technologies from class integrated into this project):
\begin{tabularx}{\linewidth}{X X}
  \toprule
  \textbf{Technology} & \textbf{Integration Status} \\
  \midrule
  \textit{[Technology 1, e.g., IPv4 routing]} & \textit{[Implemented / Partial / Not covered]} \\
  \textit{[Technology 2, e.g., DNS]} & \textit{[Implemented / Partial / Not covered]} \\
  \textit{[Technology 3, e.g., HTTP/HTTPS]} & \textit{[Implemented / Partial / Not covered]} \\
  \textit{[Technology 4, e.g., Firewall rules]} & \textit{[Implemented / Partial / Not covered]} \\
  \textit{[Technology 5, e.g., VLANs]} & \textit{[Implemented / Partial / Not covered]} \\
  \bottomrule
\end{tabularx}

\section{Flowchart}
% TODO: Visual diagram outlining the workflow, architecture, or data flow of the network.

% Include a flowchart here - use \includegraphics{flowchart.png} or draw with TikZ
% Placeholder for flowchart image:
% \begin{figure}[h!]
%   \centering
%   \includegraphics[width=0.8\textwidth]{images/flowchart.pdf}
%   \caption{Network workflow and data flow architecture}
%   \label{fig:flowchart}
% \end{figure}

% Alternatively, a TikZ diagram placeholder:
% \begin{center}
%   \begin{tikzpicture}[node distance=2cm, auto]
%     % Nodes and edges would go here
%     \node (start) {Start};
%     \node (setup) [below of start] {Network Setup};
%     \node (config) [below of setup] {Configuration};
%     \node (test) [below of config] {Testing};
%     \node (result) [below of test] {Result};
%     \node (end) [below of result] {End};
%     \path [line] (start) -- (setup);
%     \path [line] (setup) -- (config);
%     \path [line] (config) -- (test);
%     \path [line] (test) -- (result);
%     \path [line] (result) -- (end);
%   \end{tikzpicture}
% \end{center}

\section{Connections}
% TODO: Details on network topology, device interconnections, IP addressing scheme, and routing/firewall rules.

% Network Topology:
% \begin{figure}[h!]
%   \centering
%   \includegraphics[width=0.8\textwidth]{images/topology.pdf}
%   \caption{Network topology diagram}
%   \label{fig:topology}
% \end{figure}

% IP Addressing Scheme:
% \begin{tabular}{lll}
%   \toprule
%   \textbf{Device} & \textbf{IP Address} & \textbf{Subnet} \\
%   \midrule
%   \textit{[Router]} & \textit{[e.g., 192.168.1.1]} & \textit{[e.g., 255.255.255.0]} \\
%   \textit{[Server]} & \textit{[e.g., 192.168.1.10]} & \textit{[e.g., 255.255.255.0]} \\
%   \textit{[Client 1]} & \textit{[e.g., 192.168.1.20]} & \textit{[e.g., 255.255.255.0]} \\
%   \textit{[Client 2]} & \textit{[e.g., 192.168.1.30]} & \textit{[e.g., 255.255.255.0]} \\
%   \bottomrule
% \end{tabular}

% Routing and Firewall Rules:
% \begin{verbatim}
% [iptables rules / Cisco configurations / routing table entries]
% \end{verbatim}

% Example rule format:
% \begin{verbatim}
% iptables -A INPUT -p tcp --dport 80 -j ACCEPT
% iptables -A FORWARD -s 192.168.1.0/24 -d 192.168.2.0/24 -j ACCEPT
% \end{verbatim}

\section{Implementation}
% TODO: Methodology: Details on whether simulation or socket programming was used.
% TODO: Features: Description of any added checking features or click-only functionalities.

% Methodology:
% \begin{itemize}
%   \item \textbf{Approach:} \textit{[e.g., "Network simulation using Mininet" / "Socket programming in Python"]}
%   \item \textbf{Tools:} \textit{[e.g., "Mininet, Wireshark, Python 3 sockets"]}
%   \item \textbf{Implementation Steps:}
%     \begin{enumerate}
%       \item [1.] \textit{[First step, e.g., "Set up network topology"]}
%       \item [2.] \textit{[Second step, e.g., "Configure IP addressing"]}
%       \item [3.] \textit{[Third step, e.g., "Implement firewall rules"]}
%       \item [4.] \textit{[Fourth step, e.g., "Test connectivity"]}
%     \end{enumerate}
% \end{itemize}

% Features:
% \begin{itemize}
%   \item \textit{[e.g., "Real-time traffic monitoring"]}
%   \item \textit{[e.g., "Automated rule validation"]}
%   \item \textit{[e.g., "Graphical configuration interface"]}
%   \item \textit{[e.g., "Click-to-apply firewall rules"]}
% \end{itemize}

\section{Result}
% TODO: Final outputs, testing validation, and verification that the network and its rules function as expected.

% Testing Results:
% \begin{tabular}{ll}
%   \toprule
%   \textbf{Test Case} & \textbf{Result} \\
%   \midrule
%   \textit{[Connectivity test]} & \textit{[Passed / Failed]} \\
%   \textit{[Security policy test]} & \textit{[Passed / Failed]} \\
%   \textit{[Performance test]} & \textit{[Passed / Failed]} \\
%   \bottomrule
% \end{tabular}

% Verification:
% \begin{itemize}
%   \item \textit{[e.g., "Ping tests between all nodes successful"]}
%   \item \textit{[e.g., "Firewall rules blocking unauthorized access confirmed"]}
%   \item \textit{[e.g., "Bandwidth measurements within expected ranges"]}
% \end{itemize}

% Final Output Screenshots/Figures:
% \begin{figure}[h!]
%   \centering
%   \includegraphics[width=0.8\textwidth]{images/final-output.pdf}
%   \caption{Final network output / test results}
%   \label{fig:final-output}
% \end{figure}

% \begin{figure}[h!]
%   \centering
%   \includegraphics[width=0.8\textwidth]{images/packet-trace.pdf}
%   \caption{Packet trace demonstrating data flow}
%   \label{fig:packet-trace}
% \end{figure}

\end{document}